% 1_resumen.tex

\begin{abstract}
Este documento corresponde al segundo avance del proyecto de transmisión inalámbrica de audio y se presenta la propuesta de diseño del sistema junto con el protocolo de transmisión definido para su implementación.
Primero, se describe la arquitectura de dos nodos ideada: un transmisor y un receptor, cada uno basado en una Raspberry Pi Zero y un transceiver NRF24L01+PA+LNA que opera en la banda ISM de 2.4 GHz.
El nodo transmisor incorpora el micrófono digital INMP441 para la captura de audio, mientras que el nodo receptor utiliza el módulo NS4168 con parlante 3070 para la reproducción.

Asimismo, se definen los parámetros de digitalización base del sistema: audio mono a 8 kHz con resolución de 8 bits por muestra en formato PCM, con ADPCM como optimización opcional para reducir el volumen de datos transmitidos.
El protocolo de transmisión se estructura sobre el mecanismo \textit{Enhanced ShockBurst} del NRF24L01, que gestiona automáticamente la verificación CRC, las confirmaciones de recepción y las retransmisiones.
Se presentan además los diagramas de bloques de la arquitectura del sistema, el diagrama de flujo del algoritmo propuesto y la asignación de pines GPIO para cada nodo.
\end{abstract}

\begin{IEEEkeywords}
% En orden alfabético
Audio digital, banda ISM, \textit{Enhanced ShockBurst}, fragmentación de paquetes, GPIO, I2S, INMP441, NRF24L01, NS4168, Raspberry Pi Zero
\end{IEEEkeywords}
